\documentclass{article}

\usepackage{amsmath,amssymb}
\usepackage{xurl}
\usepackage[hidelinks]{hyperref}
\setlength{\emergencystretch}{2em}

% Shared commands for notes in docs/paper-gaps/.

\DeclareUnicodeCharacter{2080}{\ifmmode{}_{0}\else\textsubscript{0}\fi}
\DeclareUnicodeCharacter{2081}{\ifmmode{}_{1}\else\textsubscript{1}\fi}
\DeclareUnicodeCharacter{2082}{\ifmmode{}_{2}\else\textsubscript{2}\fi}
\DeclareUnicodeCharacter{2099}{\ifmmode{}_{n}\else\textsubscript{n}\fi}

\newcommand{\Lean}{\textsc{Lean}}
\ExplSyntaxOn
\NewDocumentCommand{\leanid}{m}{
  \ifmmode
    \textsf{\detokenize{#1}}
  \else
    \group_begin:
    \urlstyle{sf}
    \tl_set:Nn \l_tmpa_tl {#1}
    \tl_replace_all:Nnn \l_tmpa_tl {\_} {_}
    \exp_args:NV \path \l_tmpa_tl
    \group_end:
  \fi
}
\ExplSyntaxOff
\providecommand{\tr}{\operatorname{tr}}
\newcommand{\tnleanIssueBase}{https://github.com/LionSR/TNLean/issues/}
\newcommand{\tnleanPrBase}{https://github.com/LionSR/TNLean/pull/}
\newcommand{\tnleanDocBase}{https://sirui-lu.com/TNLean/docs/}
\newcommand{\ghissue}[1]{\href{\tnleanIssueBase#1}{GitHub issue \##1}}
\newcommand{\ghpr}[1]{\href{\tnleanPrBase#1}{PR \##1}}
\newcommand{\doclink}[1]{\href{\tnleanDocBase#1.html}{\path{#1}}}

% Dirac notation and the identity operator, as in blueprint/src/macros/common.tex.
% \providecommand keeps note-local definitions authoritative.
\usepackage{dsfont}
\providecommand{\ket}[1]{|#1\rangle}
\providecommand{\bra}[1]{\langle #1|}
\providecommand{\braket}[2]{\langle #1 | #2 \rangle}
\providecommand{\ketbra}[2]{|#1\rangle\!\langle #2|}
\providecommand{\Id}{\mathds{1}}
\providecommand{\one}{\mathds{1}}
\providecommand{\C}{\mathbb{C}}
\providecommand{\MN}[1]{M_{#1}(\C)}
\providecommand{\MD}{M_D(\C)}

% Verdict marker: \gapnote{<kind>}{<status>}. Kind is the policy
% classification (clarification, local-correction, scope-restriction,
% unfaithful, false-source, open-gap); status is open, wip (elimination
% underway), resolved, or historical. Machine-read by the site index;
% prints a verdict line.
\newcommand{\gapnote}[2]{\par\noindent\textbf{Verdict:} #1 (#2).\par}


\title{The Direct-Sum Input from the Older MPS Representation Paper}
\author{}
\date{2026-05-05}

\begin{document}
\maketitle
\gapnote{local-correction}{resolved}

\section{Source statement}

The finite-length direct-sum input used behind the block-injective
canonical-form route is not proved in the 2016 MPDO source. The 2016 paper
cites the older representation paper for the block-injective-after-blocking
step at lines~340--345 of Ref.~\cite{Cirac2016MPDO_arXiv}. The local source for
the older argument is \path{Papers/quant-ph_0608197/MPSarchive.tex}.

That source first introduces the injectivity condition $C1$. For matrices
$A_i\in M_D(\C)$, the condition requires a length $L_0$ for which the map
\begin{align}
    \Gamma_{L_0}(X)
     & :=
    \sum_{i_1,\ldots,i_{L_0}}
    \tr(XA_{i_1}\cdots A_{i_{L_0}})
    \ket{i_1,\ldots,i_{L_0}}
    \label{eq:direct-sum_gamma_map}
\end{align}
is injective. The source immediately identifies this condition with the
exact-length word products of length $L_0$ spanning the full matrix algebra
at lines~888--908 of Ref.~\cite{PerezGarcia2007MPSReps}.

The technical algebraic lemma states that, for every nonzero
$D\times D$ matrix $C$ and every $X\in M_D(\C)$, there are finitely many
matrices $R_i,S_i\in M_D(\C)$ such that
\begin{align}
    X
     & =\sum_i R_iCS_i.
    \label{eq:direct-sum_matrix_factorization}
\end{align}
The proof reduces $C$ to a matrix unit by elementary matrix operations and
then reconstructs $X$ from permuted copies of that matrix unit; see
lines~1333--1344 of Ref.~\cite{PerezGarcia2007MPSReps}.

The next lemma is the direct-sum statement. Its local hypotheses are important:
the tensor has $b\geq 2$ canonical blocks
$A^1,\ldots,A^b$, each block satisfies $C1$, the common length is
\begin{align}
    L_0
     & =\max_j L_0^{A^j},
    \label{eq:direct-sum_common_injectivity_length}
\end{align}
the block sizes are ordered, and the block states are pairwise different; see
lines~1321--1330 of Ref.~\cite{PerezGarcia2007MPSReps}. Under these
hypotheses, the length bound
\begin{align}
    L
     & \geq 3(b-1)(L_0+1)
    \label{eq:direct-sum_source_length_bound}
\end{align}
implies
\begin{align}
    \sum_{j=1}^b\mathcal G_L^{A^j}
     & =\bigoplus_{j=1}^b\mathcal G_L^{A^j}.
    \label{eq:direct-sum_source_image_space_directness}
\end{align}
This is the direct-sum assertion proved at lines~1346--1408 of
Ref.~\cite{PerezGarcia2007MPSReps}.

The proof proceeds by induction on the number of blocks. The base case uses
$C1$ and the matrix-factorization lemma to show that a nontrivial linear
dependence would force inclusion, then equality, of the two length-$L_0+1$ MPS
subspaces, contradicting uniqueness of the associated injective parent-space
characterization. The induction step isolates one block with a local
projector, applies the induction hypothesis to the remaining blocks, and then
uses $C1$ together with the factorization lemma to force the remaining
coefficient to vanish.

\section{MPS translation}

For the non-periodic Fundamental Theorem route, the mathematical consequence is
homogeneous and positive-length. After the canonical-form and
basis-of-normal-tensors reduction, the selected representatives should admit
one length at which their simultaneous word evaluations separate the
direct-sum blocks. In trace-dual form, the required implication is
\begin{align}
    \left[\sum_{j=1}^g\tr(\Delta_jA_j^w)=0
        \quad\text{for every word }w\text{ of length }L\right]
     & \Longrightarrow
    \Delta_j=0
    \quad\text{for every }j.
    \label{eq:direct-sum_trace_dual_separation}
\end{align}
This is the finite-length direct-sum span input needed by the BNT
pair-separation route.

The implication
(\ref{eq:direct-sum_trace_dual_separation}) is not a consequence of the empty
word and is not a cumulative-span statement. The exact length $L$ is part of
the content. Therefore the formal proof should not replace the source input by
a hypothesis that only says the cumulative pair algebra is full, nor by a
theorem that uses the length-zero identity as padding.

\section{Formal boundary}

The matrix-factorization content of Lemma~\texttt{lem1} from the source paper
is now formalized in two steps.\footnote{%
    \leanid{Matrix.span_range_mul_nonzero_mul_eq_top} for the abstract
    matrix-algebra form;
    \leanid{MPSTensor.span_range_evalWord_mul_nonzero_mul_evalWord_eq_top}
    for the word-product translation.}
The first says that, for every nonzero matrix $C$,
\begin{align}
    \operatorname{span}
    \{RCS:R,S\in M_D(\C)\}
     & =M_D(\C).
    \label{eq:direct-sum_nonzero_matrix_product_span}
\end{align}
The second translates this assertion into the relevant word-product form. If a
tensor is $L$-block injective and $X\ne 0$, then
\begin{align}
    \operatorname{span}
    \{A_uXA_v:|u|=|v|=L\}
     & =M_D(\C).
    \label{eq:direct-sum_word_product_factorization_span}
\end{align}

The next part of the two-block direct-sum argument is also formalized. In the
size-ordering convention $D_B\leq D_A$, the three-block trace relation gives
\begin{align}
    \mathcal G_L^A
     & \subseteq\mathcal G_L^B.
    \label{eq:direct-sum_ground_space_inclusion}
\end{align}
Block injectivity identifies the dimensions as
\begin{align}
    \dim\mathcal G_L^A
     & =D_A^2,
    \label{eq:direct-sum_A_ground_space_dimension} \\
    \dim\mathcal G_L^B
     & =D_B^2.
    \label{eq:direct-sum_B_ground_space_dimension}
\end{align}
The inclusion
(\ref{eq:direct-sum_ground_space_inclusion}), the dimension identities
(\ref{eq:direct-sum_A_ground_space_dimension}) and
(\ref{eq:direct-sum_B_ground_space_dimension}), and the size ordering together
force
\begin{align}
    D_A
     & =D_B,
    \label{eq:direct-sum_forced_equal_bond_dimension} \\
    \mathcal G_L^A
     & =\mathcal G_L^B.
    \label{eq:direct-sum_forced_equal_ground_spaces}
\end{align}
This implication is formalized by
\leanid{MPSTensor.bondDim_eq_and_groundSpace_eq_of_three_block_trace_relation_left_of_dim_ge}.

The strict-size branch states that, if $D_B<D_A$, no nonzero coefficient on
the larger block can satisfy the homogeneous three-block trace
relation.\footnote{%
    \leanid{MPSTensor.not_three_block_trace_relation_left_of_isNBlkInjective_of_dim_gt};
    \leanid{MPSTensor.pairTraceSeparatingAt_threeBlock_of_isNBlkInjective_of_dim_gt}
    for the trace-separation form;
    \leanid{MPSTensor.pairTraceSeparatingAt_threeBlock_of_isNBlkInjective_of_dim_ne}
    for the convention-free version.}
The conditional directness result is also formalized: an external proof of
\begin{align}
    \neg\left(D_A=D_B
    \mathbin{\wedge}
    \mathcal G_L^A=\mathcal G_L^B\right)
    \label{eq:direct-sum_excluded_equal_size_collapse}
\end{align}
implies
\begin{align}
    \mathcal G_{3L}^A\cap\mathcal G_{3L}^B
     & =\{0\}.
    \label{eq:direct-sum_three_block_zero_intersection}
\end{align}
The corresponding declaration is
\leanid{MPSTensor.groundSpace_inf_eq_bot_of_not_bondDim_eq_and_groundSpace_eq_of_dim_ge}.

The equal-size uniqueness branch is isolated as well. For injective tensors,
equality of the length-$L$ image spaces forces equality of the periodic MPV
lines. A supplied distinct-MPV-line hypothesis therefore rules out the
equal-size collapse.\footnote{%
    \leanid{MPSTensor.not_bondDim_eq_and_groundSpace_eq_of_mpvSubmodule_ne}.}
The analytic source of distinct-MPV witnesses is formalized for irreducible
trace-preserving blocks with normalized self-overlaps: if two blocks are not
gauge-phase equivalent, then at arbitrarily large chain lengths their MPV
states are not pointwise proportional.\footnote{%
    \leanid{MPSTensor.exists_ge_not_forall_mpv_eq_mul_of_not_gaugePhaseEquiv_of_irreducible_TP};
    \leanid{MPSTensor.exists_ge_not_forall_mpv_eq_mul_of_blocksNotGaugePhaseEquiv_of_irreducible_TP}
    for the BNT-family form.}

Combining the uniqueness argument with the finite-dimensional image-space step
gives the following implication: if some chain length
\begin{align}
    N
     & \geq\max\{2,L\}
    \label{eq:direct-sum_nonproportionality_witness_length}
\end{align}
witnesses nonproportionality, then
(\ref{eq:direct-sum_three_block_zero_intersection})
holds.\footnote{%
    \leanid{MPSTensor.groundSpace_inf_eq_bot_of_exists_not_forall_mpv_eq_mul_of_dim_ge};
    \leanid{MPSTensor.groundSpace_inf_eq_bot_of_blocksNotGaugePhaseEquiv_same_dim_of_dim_ge}
    for the BNT-separation packaging.}
Zero intersection of image spaces at a fixed length, together with injectivity
of the boundary-to-chain maps at that length, implies homogeneous pair trace
separation.\footnote{%
    \leanid{MPSTensor.pairTraceSeparatingAt_of_groundSpace_inf_eq_bot_of_isNBlkInjective};
    and the combined forms
    \leanid{MPSTensor.pairTraceSeparatingAt_threeBlock_of_exists_not_forall_mpv_eq_mul_of_dim_ge},
    \leanid{MPSTensor.pairTraceSeparatingAt_threeBlock_of_blocksNotGaugePhaseEquiv_same_dim_of_dim_ge}.}
The two branches are combined for all ordered pairs of distinct blocks of a
BNT family at the common length $3L$, and the result is fed into the selector
and product-span construction.\footnote{%
    \leanid{MPSTensor.forall_pairTraceSeparatingAt_threeBlock_of_blocksNotGaugePhaseEquiv};
    \leanid{MPSTensor.exists_forall_pairTraceSeparatingAt_threeBlock_of_blocksNotGaugePhaseEquiv}.}

The BNT conclusion is stated directly from explicit separated-block
hypotheses, with the required finite-length injectivity hypotheses visible.
These hypotheses may begin at an arbitrary positive common length $L_0$;
positive-length propagation supplies the longer lengths used in the proof.
The selector datum and the resulting finite direct-sum span are theorems rather
than identity-padding assumptions.\footnote{%
    \leanid{MPSTensor.hasPairBlockSeparatingWords_threeBlock_of_blocksNotGaugePhaseEquiv};
    \leanid{MPSTensor.propBlockInjective_of_common_blockInjective_of_pairBlockSeparatingWords};
    \leanid{MPSTensor.wordTupleSpanTop_of_blocksNotGaugePhaseEquiv_directSum_selectors_c1};
    all in
    \path{TNLean/MPS/BNT/DirectSumSelectors.lean}.}

The sharp direct-sum interpolation does not concatenate a separate injective
prefix. Full pair span at a common length $S$ realizes an arbitrary value
$(X,0)$ on a chosen block and one anchor block. Multiplication by $b-2$
identity--zero selectors removes the remaining blocks. Hence the full
simultaneous span holds at length
\begin{align}
    (b-1)S.
    \label{eq:direct-sum_selector_interpolation_length}
\end{align}
For
\begin{align}
    S
     & =3(L_0+1),
    \label{eq:direct-sum_three_block_pair_length}
\end{align}
the resulting length is
\begin{align}
    (b-1)S
     & =3(b-1)(L_0+1),
    \label{eq:direct-sum_recovered_source_length}
\end{align}
which is exactly the source length in
(\ref{eq:direct-sum_source_length_bound}).\footnote{%
    \leanid{MPSTensor.wordTupleSpanTop_mul_pred_of_forall_pairWordTupleSpanTop};
    \leanid{MPSTensor.wordTupleSpanTop_threeBlock_mul_pred_of_blocksNotGaugePhaseEquiv_c1}.}

The downstream consumers of the direct-sum input are threefold. The
exact-length span of a direct-sum block family and its two-block trace-dual
form are expressed as predicates used throughout the pair route.\footnote{%
    \leanid{MPSTensor.WordTupleSpanTop}, \leanid{MPSTensor.PairWordTupleSpanTop},
    \leanid{MPSTensor.PairTraceSeparatingAt}.}
Theorems in the period-window module turn a positive homogeneous window of the
pair-span predicate into homogeneous trace separation.\footnote{%
    \path{TNLean/MPS/MPDO/BiCFDerivation/PairHomogenization.lean}.}
The BNT direct-sum theorems first turn pairwise non-equivalence into homogeneous
three-block trace separation, and then construct selectors and a common
exact-length simultaneous span.\footnote{%
    \leanid{MPSTensor.forall_pairTraceSeparatingAt_threeBlock_of_blocksNotGaugePhaseEquiv};
    \leanid{MPSTensor.hasPairBlockSeparatingWords_threeBlock_of_blocksNotGaugePhaseEquiv};
    \leanid{MPSTensor.wordTupleSpanTop_of_blocksNotGaugePhaseEquiv_directSum_selectors_c1}
    in \path{TNLean/MPS/BNT/DirectSumSelectors.lean}; and
    \leanid{MPSTensor.wordTupleSpanTop_mul_pred_of_forall_pairWordTupleSpanTop},
    \leanid{MPSTensor.hasBlockSelectorWords_of_pairBlockSeparatingWords}
    in \path{TNLean/MPS/MPDO/BiCFDerivation/Selectors.lean}.}

\section{The distinctness hypothesis}

The formal results above assume that no two distinct blocks are related by a
gauge and a phase: there are no indices $j\ne k$ with equal block dimensions,
invertible matrix $X$ and scalar $\zeta\ne0$ such that
\begin{align}
    A^k_i
     & =\zeta XA^j_iX^{-1}
    \quad\text{for every }i.
    \label{eq:direct-sum_gauge_phase_relation}
\end{align}
The source states its hypothesis on the block states
\begin{align}
    \ket{\phi_{A^j}}
     & =\lambda_j^N\sum_{i_1,\ldots,i_N}
    \tr(A^j_{i_1}\cdots A^j_{i_N})\ket{i_1\cdots i_N}
    \label{eq:direct-sum_block_state}
\end{align}
of the $N$-site ring: without loss of generality, these states are pairwise
different (lines~1326--1330 of Ref.~\cite{PerezGarcia2007MPSReps}).

The intended meaning of ``different'' is ``not proportional'', that is,
different rays. Three passages of the source fix this reading.
\begin{enumerate}
    \item The source writes equality of states for equality up to a scalar.
          The uniqueness proof for periodic chains ends with
          $X=\lambda I$ ``and $\ket\phi=\ket\psi$'' (line~1289), and the
          open-boundary proof fixes the remaining scalar by normalization
          (lines~1233--1236).
    \item The base case of the direct-sum lemma ends by observing that both
          $\ket{\phi_{A^1}}$ and $\ket{\phi_{A^2}}$ would be the only ground
          state of one parent Hamiltonian, ``which is the desired
          contradiction'' (lines~1373--1377). That uniqueness theorem determines
          a ground state only up to a scalar, so this is a contradiction exactly
          when the two block states are not proportional.
    \item The phrase ``without loss of generality'' refers to merging blocks
          whose states are proportional. Two such terms add to a multiple of
          one block state, which is again of the form
          (\ref{eq:direct-sum_block_state}) after rescaling and rephasing the
          block; when the two terms cancel, both are dropped. Merging only
          equal vectors would leave the proportional pairs for which the
          contradiction above fails.
\end{enumerate}
The source attaches no separate qualifier to block states at comparable
statements. Where it relates matrices up to a phase it says so
(lines~1013--1014), and the proof of the corollary on the W state counts
``the number of different blocks'' in the same sense (line~2189).

At the lengths of the source the two hypotheses coincide. If
(\ref{eq:direct-sum_gauge_phase_relation}) holds, cyclicity of the trace gives,
for every length $N$,
\begin{align}
    \ket{\phi_{A^k}}
     & =\left(\frac{\lambda_k\zeta}{\lambda_j}\right)^N\ket{\phi_{A^j}},
    \label{eq:direct-sum_gauge_phase_proportional}
\end{align}
so pairwise nonproportional block states exclude
(\ref{eq:direct-sum_gauge_phase_relation}). Conversely, under gauge-phase
inequivalence, condition $C1$ in each block and
$N\ge\max\{L_0,3(b-1)(L_0+1)\}$, the block vectors at length $N$ are linearly
independent, hence pairwise nonproportional.\footnote{%
    \leanid{MPSTensor.PGVWC07CanonicalFormData.wordTupleSpanTop_of_ge} and
    \leanid{MPSTensor.linearIndependent_mpvState_of_wordTupleSpanTop}.}
Gauge-phase inequivalence is therefore the length-independent form of the
source's distinctness hypothesis.

\textbf{Local fix (projective distinctness of the blocks).} The library reads
``pairwise different block states'' as pairwise nonproportional block states
and states it as gauge-phase inequivalence of distinct blocks. It is an
explicit premise of the theorems that use it, not a field of the canonical-form
data. This placement follows the source. Theorem~4 of
Ref.~\cite{PerezGarcia2007MPSReps} (lines~742--763) imposes no distinctness,
and nothing in its proof merges blocks related by a gauge and a phase. The
distinctness enters only through the standing hypotheses of the degeneracy
section, lines~1321--1330, and the formal theorems carry it exactly where the
source does: at the direct-sum lemma and at the degeneracy theorems that use
it. The literal reading, in which two block states may be unequal but
proportional, is not one the authors intend and receives no separate formal
statement.

\section{Conclusion}

Under the distinctness hypothesis, the arbitrary-$L_0$ direct-sum interpolation
needed for the BNT specialization is formalized at the exact source length
\begin{align}
    3(b-1)(L_0+1).
    \label{eq:direct-sum_exact_formalized_length}
\end{align}
The resulting span is homogeneous and positive-length, not cumulative. The
length-zero identity is only the neutral base of a finite product over no
remaining selectors. The case $b=1$ is treated separately by ordinary block
injectivity.

For a BNT canonical form of total bond dimension $D$, the common quantum
Wielandt length is
\begin{align}
    L_0
     & =D^4.
    \label{eq:direct-sum_quantum_wielandt_length}
\end{align}
Every representative has positive dimension and occurs at least once, so
$b\leq D$. The formal theorem consequently gives
\begin{align}
    3(b-1)(D^4+1)
     & \leq 3(D-1)(D^4+1)
    \label{eq:direct-sum_basis_count_dimension_bound} \\
     & \leq 3D^5.
    \label{eq:direct-sum_total_dimension_length_bound}
\end{align}
This proves the sharp block-separation input cited by the 2016 MPDO source
with the source's distinctness hypothesis in the form of gauge-phase
inequivalence.\footnote{%
    \leanid{MPSTensor.SectorDecomposition.basisCount_le_totalDim};
    \leanid{MPSTensor.IsBNTCanonicalForm.exists_basis_wordTupleSpanTop_le_three_totalDim_pow_five}.}

This conclusion concerns the finite-length direct-sum span. Other statements
in the Fundamental Theorem that compare two canonical-form tensors still
require their own reblocking and sector-data arguments; those are logically
separate from the direct-sum input proved here.

\section{Arbitrary positive interactions}

The same distinctness hypothesis enters Theorem~11 of
Ref.~\cite{PerezGarcia2007MPSReps}, lines~1407--1422, which concerns an arbitrary
positive local interaction, rather than the canonical parent projection.
Let $h$ act on $L$ consecutive sites, and let $H=\sum_i\tau^i(h)$ on the
$N$-site ring. If the complementary word tuples span the product matrix
algebra at length $N-L$, then
\begin{align}
    H\sum_j c_j V^{(N)}(A_j)=0,\quad c_j\ne0
    \quad\Longrightarrow\quad
    HV^{(N)}(A_j)=0\quad\text{for every }j.
    \label{eq:direct_sum_arbitrary_interaction_blocks}
\end{align}
Indeed, positivity first makes each local translate annihilate the weighted
sum. At any cyclic cut, the complementary word is the boundary matrix for the
active block. Applying $h$ to that boundary map gives a linear map on the
product of block matrix algebras. Its vanishing on the complementary word
family extends to the whole product algebra; a boundary supported in a single
block isolates that block's local annihilation equation. Summing the local
equations gives (\ref{eq:direct_sum_arbitrary_interaction_blocks}).

Canonical data, common injectivity at $L_0>0$, and gauge-phase inequivalence
derive the complementary span when $b\ge2$ and $N\ge3(b-1)(L_0+1)+L$. The
resulting theorem has the hypotheses of Theorem~11, with the distinctness of the
block states read as above, and its exact volume bound.\footnote{%
    \leanid{MPSTensor.periodicInteractionHamiltonianES_mpvState_eq_zero_of_wordTupleSpanTop};
    \leanid{MPSTensor.PGVWC07CanonicalFormData.periodicInteractionHamiltonianES_blocks_eq_zero}.}

\section{Verdict}

\textbf{Verdict: local correction, resolved.} The source's pairwise different
block states are pairwise nonproportional vectors, as its own proof of the
direct-sum lemma requires. At the source lengths this is equivalent to
gauge-phase inequivalence of distinct blocks, the explicit premise of the formal
direct-sum span and of the degeneracy theorem for arbitrary positive
interactions. The exact source length is proved under this hypothesis.

\bibliographystyle{alpha}
\bibliography{references}

\end{document}
